\documentclass[12pt,oneside]{amsart}

%\usepackage[T1]{fontenc}
%\usepackage[utf8]{inputenc}
\usepackage[german]{babel}
\usepackage{geometry}                % See geometry.pdf to learn the layout options. There are lots.
\geometry{a4paper}b                   % ... or a4paper or a5paper or ... 
\usepackage{color}
\usepackage[table]{xcolor}
\usepackage{graphicx}
\usepackage[mtphrb,mtpcal,mtpfrak]{mtpro2}
\usepackage{amsmath}
\usepackage{amsthm}
\usepackage{enumerate}
\usepackage{tikz}
\usetikzlibrary{calc}
\usetikzlibrary{arrows.meta}
\usepackage[shortalphabetic]{amsrefs}
\usepackage{wrapfig}
\usepackage[textsize=tiny]{todonotes}
\usepackage[many]{tcolorbox}
\usepackage{pgfplots}
\usepackage{booktabs}
\usepgfplotslibrary{patchplots}
\usepgfplotslibrary{fillbetween}

%\usepackage[no-math]{fontspec}
%\setsansfont[Scale=0.92]{Helvetica Neue LT Std 65 Medium}
%\usepackage[scaled=0.92]{helvet}
%\usepackage{times}



\usepackage{mathspec}
\setmathrm{Times LT Std}
\setmainfont[BoldFont=Times LT Std Bold, ItalicFont=Times LT Std Italic, BoldItalicFont=Times LT Std Bold Italic, Ligatures=TeX]{Times LT Std}

%\setmainfont{Baskerville URW}

%\setsansfont[Scale=0.92,BoldFont=Helvetica Neue LT Std 65 Medium,
%    ItalicFont=Helvetica Neue LT Std 56 Italic,
%    BoldItalicFont=Helvetica Neue LT Std 66 Medium Italic,
%    Ligatures=TeX]{Helvetica Neue LT Std 55 Roman}
\setsansfont[BoldFont=Whitney-Semibold.otf,ItalicFont=Whitney-MediumItalic.otf,Ligatures=TeX]{Whitney-Medium.otf}
\setboldmathrm[BoldFont={Times LT Std Bold}]{Times LT Std Bold}

\newcommand{\red}[1]{\textcolor{red}{#1}}
\definecolor{light-gray}{gray}{0.85}

%%% here we define our palette

\definecolor{spacecadet}{HTML}{0D284C}
\definecolor{munsell}{HTML}{008FA8}
\definecolor{banana}{HTML}{FFD932}
\definecolor{cgblue}{HTML}{007CA5}
\definecolor{isabelline}{HTML}{EAEDEA}

%%%

\makeatletter
\def\@secnumfont{\bfseries}
\def\part{\@startsection{part}{0}%
  \z@{\linespacing\@plus\linespacing}{.5\linespacing}%
  {\normalfont\bfseries\raggedright}}
\def\specialsection{\@startsection{section}{1}%
  \z@{\linespacing\@plus\linespacing}{.5\linespacing}%
  {\color{cgblue}\normalfont\centering}}
\def\section{\@startsection{section}{1}%
  \z@{.7\linespacing\@plus\linespacing}{.5\linespacing}%
  {\color{cgblue}\normalfont\sf\bfseries\large\centering}}
\def\subsection{\@startsection{subsection}{2}%
  \z@{.5\linespacing\@plus.7\linespacing}{-.5em}%
  {\color{cgblue}\normalfont\sf\bfseries}}
\def\subsubsection{\@startsection{subsubsection}{3}%
  \z@{.5\linespacing\@plus.7\linespacing}{-.5em}%
  {\color{cgblue}\normalfont\sf\itshape}}
\def\paragraph{\@startsection{paragraph}{4}%
  \z@\z@{-\fontdimen2\font}%
  \color{cgblue}\sf\normalfont}
\def\subparagraph{\@startsection{subparagraph}{5}%
  \z@\z@{-\fontdimen2\font}%
  \color{cgblue}\sf\normalfont}
\makeatother

\DeclareMathOperator{\supp}{supp}
\DeclareMathOperator{\conf}{conf}
\DeclareMathOperator{\tr}{tr}
\DeclareMathOperator{\im}{im}
\DeclareMathOperator{\id}{id}
\DeclareMathOperator{\vspan}{span}
\DeclareMathOperator{\res}{res}
\DeclareMathOperator{\real}{Re}
\DeclareMathOperator{\imag}{Im}
\DeclareMathOperator{\diff}{diff}
\DeclareMathOperator{\homeo}{Homeo}
\DeclareMathOperator{\Mor}{Mor}
\DeclareMathOperator{\Ob}{Ob}
\DeclareMathOperator{\target}{Target}
\DeclareMathOperator{\source}{Source}
\DeclareMathOperator{\aut}{Aut}
\DeclareMathOperator{\Sch}{Sch}

\theoremstyle{plain}% default
\newtheorem{theorem}{Theorem}[section]
\newtheorem{lemma}[theorem]{Lemma}
\newtheorem{proposition}[theorem]{Proposition}
\newtheorem*{corollary}{Corollary}

\theoremstyle{definition}
\newtheorem{definition}{\color{cgblue}Definition}[section]
\newtheorem{prototypedefinition}{Prototype Definition}[section]
\newtheorem{physicalintuition}{Physical intuition}[section]
\newtheorem{conjecture}{Conjecture}[section]
\newtheorem{beispiel}{\color{cgblue}Beispiel}[section]
\newtheorem{übung}{\color{cgblue}Übung}[section]
\newtheorem{exercise}{\color{cgblue}Exercise}[section]

\theoremstyle{remark}
\newtheorem*{remark}{Remark}
\newtheorem*{note}{Note}

\newenvironment{lösung}
  {\begin{proof}[\color{cgblue}Lösung]\color{spacecadet}}
  {\end{proof}}

\def\eqbox#1{\tcboxmath[colback=banana,arc=0cm,frame hidden,enhanced]{#1}}
\def\figbox{\begin{center}\tcboxmath[colback=isabelline,frame hidden,enhanced,sharp corners]{\parbox[c][3cm][c]{10cm}{\ }}\end{center}}

\title[Quantum computing and quantum logic with trapped ions]{Quantum computing and quantum logic with trapped ions: Homework 10, due Jan.\ 13th, 2023}
\author{Tobias J.\ Osborne}
\date{\today}                                           % Activate to display a given date or no date

\begin{document}

\maketitle
Groups of two students can submit together. Calculations with computer algebra systems are acceptable if well documented (please also upload exported PDF). There will be a submission folder on Stud.IP. Please upload your solutions electronically and clearly state whom you are submiting for on the submission (name and student registration number)!

\setcounter{section}{1}

\begin{exercise}[2 marks]
  We argued in the lecture that a standard query, which maps $|j,b\rangle\mapsto |j,b\oplus x_j\rangle$ (where
$j\in \{0,\ldots , N − 1\}$ and $b \in \{0, 1\}$), can be used to implement a phase-query to $x$, i.e., one of
the type $|j\rangle \mapsto (-1)^{x_j}|j\rangle$ (this is an uncontrolled phase-query). Show that a standard query can be implemented using controlled phase-queries to $x$ (which map $|c,j\rangle \mapsto (-1)^{cx_j}|c,j\rangle$ so the phase is added only if the control bit is $c = 1$), and possibly some auxiliary qubits and other gates.
\end{exercise}

\begin{exercise}[5 marks]
  Suppose we have a 2-bit input $x = x_0 x_1$ and a phase-query that maps
  \begin{equation}
    O_{x,\pm}:|b\rangle \mapsto (-1)^{x_b}|b\rangle, \quad \text{for $b\in \{0,1\}$.}
  \end{equation}
  Suppose we run the $1$-qubit circuit $HO_{x,\pm}H$ on the initial state $|0\rangle$ and then measure (in the computational basis). What is the probability distribution on the output bit, as a function of $x$? Now suppose the query leaves some workspace in a second qubit, which is initially $|0\rangle$:
  \begin{equation}
    O'_{x,\pm}:|b,0\rangle \mapsto (-1)^{x_b}|b,b\rangle, \quad \text{for $b\in \{0,1\}$.}
  \end{equation}
  Suppose we just ignore the workspace and run the algorithm of (a) on the first qubit with $O'_{x,\pm}$ instead of $O_{x,\pm}$ (and $H\otimes \mathbb{I}$ instead of $H$, and initial state $|00\rangle$). What is now the probability distribution on the output bit (i.e., if we measure the first of the two bits)?
\end{exercise}

\begin{exercise}[5 marks]
  Suppose $n = 2$, and $x = x_{00}x_{01}x_{10}x_{11} = 0001$. Give the specific initial state, three intermediate states, and final state in Grover’s algorithm, for $k = 1$ iterations. What is the success probability? Give the final state after k = 2 iterations. What is now the success probability?
\end{exercise}

\begin{exercise}[3 marks]
  Show that if the number of solutions is $t = N/4$, then Grover’s algorithm always finds a solution with certainty after just one query. How many queries would a classical algorithm need to find a solution with certainty if $t = N/4$? And if we allow the classical algorithm error probability $1/10$?
\end{exercise}

\end{document}
